\section{Discussion}
\subsubsection{Measurement.} The confirmatory result that survives multiplicity correction is H3: an English keyword list, the most common operationalisation, identified about one third of the AI-related notices estimated from the design-based sample. Adding core terms in the EU languages doubled recall without loss of precision, so most of the gain came from handling language rather than from a more complex model. This is consistent with the multilingual composition of TED \cite{ahmia2018} and with the transparency--recall trade-off discussed for dictionary measures \cite{gentzkow2019}. The language-specific gap predicted by H1 appeared only when contested labels were removed; under the primary reference it was masked by the uncertainty from a few heavily weighted positives.

\subsubsection{Why the supervised classifiers failed.} Contrary to H2, the embedding and TF-IDF classifiers underperformed the multilingual dictionary. Two features of the design explain this post hoc. First, training with inverse-probability weights gave the 450 stratum-D notices weights about 200 times larger than the dictionary strata, which pushed predicted probabilities towards zero; unweighted training raised F1 from 0.06 to 0.35. Second, with 364 positives, most of which contain explicit AI terms, a linear classifier on general-purpose sentence embeddings learned little beyond the dictionary signal. These results do not show that supervised measurement is ineffective in general; they show that, at this label budget, weighting choices that are appropriate for estimation can be harmful for training, an implementation risk that is easy to overlook.

\subsubsection{Downstream inference.} The naive multilingual-dictionary regression produced a statistically significant negative association between AI-relatedness and single bidding, whereas design-based estimates were too imprecise to confirm or reject any association. The precision of DSL depends on how well $\hat{g}$ predicts the rare positives in the low-probability stratum \cite{egami2023,audinet2025}; here, four contested labels determined much of the variance. The practical lesson is that a significant coefficient obtained with a keyword indicator should not be read as evidence about AI contracts when the measurement error has not been quantified \cite{battaglia2025}. For procurement monitoring, the results favour a multilingual dictionary for screening, combined with a labelled probability sample that devotes more of the budget to the unflagged mass, where most of the uncertainty lies.

\section{Limitations}
The reference labels come from an LLM codebook annotation rather than from expert coders, and agreement was only moderate for English notices; the construct is therefore ``AI-related according to this codebook and annotator'', and the four contested stratum-D positives show that conclusions about prevalence depend on this choice. The frame covers only notices in two CPV divisions published in TED, i.e., mostly above EU thresholds, so AI purchases classified elsewhere or below thresholds are not observed. The regression excludes multi-lot notices and notices without a valid tender count, which may introduce selection, and the controls omit buyer experience and supplier history, so omitted-variable bias is possible; all associations are non-causal. The dictionaries were written by the authors and may contain translation errors, and the stratification used M1, M2, and M4 scores, which makes the precision of estimates for other methods dependent on this design. Finally, a single small encoder and linear classifiers were tested; fine-tuned or larger models may perform differently.

\section{Conclusion}
We evaluated how AI-related public procurement notices are identified in 54,160 multilingual TED award notices, using a pre-specified stratified reference sample and design-based estimation. A multilingual dictionary roughly doubled the recall of the common English keyword approach, the English approach captured about one third of the estimated AI notices, and an apparent negative association between AI-relatedness and single bidding under naive measurement could not be confirmed after accounting for measurement error. Future work should use expert-adjudicated labels, allocate more labels to the unflagged stratum, test fine-tuned multilingual classifiers, and extend the analysis to national systems such as the Korean KONEPS.
